% Clase 1 — Búsqueda Adversaria y Juegos (parte 1)
% Lectura: AIMA Cap. 6
\section{Búsqueda Adversaria y Juegos I}

\subsection{Juegos como Problemas de Búsqueda}

Un \textbf{juego} de dos jugadores, suma cero, de información perfecta se define como:
\begin{itemize}
    \item $S_0$: estado inicial
    \item $\text{PLAYER}(s)$: qué jugador mueve en estado $s$
    \item $\text{ACTIONS}(s)$: movimientos legales
    \item $\text{RESULT}(s, a)$: estado resultante
    \item $\text{IS-TERMINAL}(s)$: test de fin de juego
    \item $\text{UTILITY}(s, p)$: valor numérico del estado terminal para el jugador $p$
\end{itemize}

\subsection{Minimax}

MAX maximiza, MIN minimiza la utilidad.

$$\text{MINIMAX}(s) = \begin{cases}
\text{UTILITY}(s, MAX) & \text{si IS-TERMINAL}(s) \\
\max_{a} \text{MINIMAX}(\text{RESULT}(s,a)) & \text{si PLAYER}(s) = MAX \\
\min_{a} \text{MINIMAX}(\text{RESULT}(s,a)) & \text{si PLAYER}(s) = MIN
\end{cases}$$

\begin{itemize}
    \item Completo (en árbol finito) y óptimo contra un oponente óptimo.
    \item Tiempo: $\bigO{b^m}$ \quad Espacio: $\bigO{bm}$
\end{itemize}
